
\section*{Глава 3. Разработка веб-приложения}
\addcontentsline{toc}{section}{Глава 3. Разработка веб-приложения}
\stepcounter{section}

\subsection{Разработка пользовательского интерфейса и реализация базовой структуры страниц}

В соответствии с задачей~3.1 была реализована серверная шаблонизация на основе
движка Jinja2, встроенного в фреймворк FastAPI. Все страницы
приложения строятся через наследование от единого базового шаблона
base.html.

Навигационное меню формируется динамически в зависимости от роли вошедшего
пользователя: администратор видит ссылку на панель управления, художественный
руководитель~--- дополнительный пункт «Отчеты», остальные роли получают только
доступ к разделу мероприятий. Пример логики условного отображения меню
приведён в листинге~\ref{lst:nav_template}.

\begin{lstlisting}[caption={Фрагмент шаблона навигационного меню}, label={lst:nav_template}]
<nav class="main-nav">
  {% if current_user.role == 'ADMIN' %}
    <a href="/admin">Панель администратора</a>
  {% else %}
    <a href="/events">Мероприятия</a>
    {% if current_user.role == 'ART_DIRECTOR' %}
      <a href="/reports">Отчеты</a>
    {% endif %}
  {% endif %}
  <a href="/logout">Выйти</a>
</nav>
\end{lstlisting}

Корневой маршрут / перенаправляет пользователя на страницу
мероприятий с помощью RedirectResponse, что позволяет избежать
пустой стартовой страницы. Статические файлы (шрифты, CSS) обслуживаются
через механизм StaticFiles фреймворка FastAPI по префиксу /static.

Для единообразия визуального оформления в приложении используется
фирменный шрифт «George» в четырёх начертаниях (Regular, SemiBold, Bold,
Italic). Цветовая палитра и шрифт соответствуют фирменному стилю ДК Подмосковья ~\cite{brandbook}.
На страницах со списками мероприятий реализована постраничная навигация
(пагинация): число записей на странице фиксировано, а ссылки на соседние
страницы формируются шаблонным тегом с передачей текущих параметров
фильтрации (см. листинг~\ref{lst:pagination}).

\begin{lstlisting}[caption={Фрагмент шаблона пагинации}, label={lst:pagination}]
{% if total_pages > 1 %}
<nav class="pagination">
  {% if page > 1 %}
    <a href="?search={{ search }}&status={{ status_filter }}
             &page={{ page - 1 }}" class="page-btn">&laquo;</a>
  {% endif %}
  {% for p in range(1, total_pages + 1) %}
    <a href="?search={{ search }}&status={{ status_filter }}&page={{ p }}"
       class="page-btn {% if p == page %}active{% endif %}">{{ p }}</a>
  {% endfor %}
  {% if page < total_pages %}
    <a href="?search={{ search }}&status={{ status_filter }}
             &page={{ page + 1 }}" class="page-btn">&raquo;</a>
  {% endif %}
</nav>
{% endif %}
\end{lstlisting}
\newpage
\subsection{Реализация модуля аутентификации пользователей и разграничения прав доступа}

Согласно задаче~3.2, в системе реализована многоуровневая модель доступа на
основе ролей (RBAC). Предусмотрены четыре роли:
\begin{itemize}
    \item ADMIN --- системный администратор, управляет учётными
          записями пользователей;
    \item ART\_DIRECTOR --- художественный руководитель, имеет
          полный доступ к мероприятиям и отчётам;
    \item DEPARTMENT\_HEAD --- руководитель отдела, создаёт
          мероприятия и назначает ответственных;
    \item EMPLOYEE --- сотрудник, редактирует мероприятия,
          за которые назначен ответственным.
\end{itemize}

Для обеспечения безопасности применяется протокол JWT (JSON Web Token) в связке с библиотекой python-jose. При успешной аутентификации сервер генерирует токен, содержащий идентификатор пользователя в поле sub и временную метку истечения exp.
Токен подписывается секретным ключом по алгоритму HS256.
(см.\ листинг~\ref{lst:jwt_create}).

\begin{lstlisting}[caption={Функция генерации JWT-токена}, label={lst:jwt_create}]
def create_access_token(
    subject: str,
    expires_delta: timedelta | None = None,
) -> str:
    expire = datetime.now(UTC) + (
        expires_delta if expires_delta is not None
        else timedelta(minutes=settings.jwt_access_token_expire_minutes)
    )
    payload: dict[str, Any] = {"sub": subject, "exp": expire}
    return jwt.encode(
        payload,
        settings.jwt_secret_key,
        algorithm=settings.jwt_algorithm,
    )
\end{lstlisting}

Пароли хранятся в базе данных в хешированном виде с использованием
библиотеки bcrypt.

Для верификации входящих запросов зависимость
get\_current\_user (модуль app/api/dependencies.py)
извлекает токен из куки access\_token, декодирует его,
получает из поля sub идентификатор пользователя и загружает
соответствующую запись из базы данных. Если токен недействителен
или пользователь не найден, возвращается HTTP-ответ
401 Unauthorized.

Разграничение прав доступа к операциям над мероприятиями реализовано в
модуле app/core/permissions.py через набор чистых функций
(листинг~\ref{lst:permissions}).

\begin{lstlisting}[caption={Функции проверки прав доступа}, label={lst:permissions}]
def can_create_event(user: User) -> bool:
    return user.role in (Role.DEPARTMENT_HEAD, Role.ART_DIRECTOR)

def can_edit_event(user: User, event: Event) -> bool:
    if user.role == Role.ART_DIRECTOR:
        return True
    if user.role in (Role.EMPLOYEE, Role.DEPARTMENT_HEAD):
        return event.responsible_id == user.id
    return False

def can_delete_event(user: User) -> bool:
    return user.role == Role.ART_DIRECTOR

def can_assign_responsible(user: User) -> bool:
    return user.role in (Role.DEPARTMENT_HEAD, Role.ART_DIRECTOR)
\end{lstlisting}

Функции используются непосредственно в роутерах: при нарушении
правил доступа возвращается ответ 403 Forbidden с текстом
«Недостаточно прав».
\newpage
\subsection{Разработка CRUD-интерфейсов для управления мероприятиями и документами}

В соответствии с задачей~3.3 разработаны полные CRUD-интерфейсы для
двух ключевых сущностей: мероприятий и прикреплённых
документов.

REST API для работы с мероприятиями реализован в модуле
app/api/events.py. Поддерживаются следующие операции:

\begin{itemize}
    \item GET /api/v1/events~--- получить список всех
          мероприятий с возможностью фильтрации по статусу и поиску
          по названию;
    \item GET /api/v1/events/\{id\}~--- получить одно
          мероприятие по идентификатору;
    \item POST /api/v1/events~--- создать новое мероприятие
          (только для ролей DEPARTMENT\_HEAD и
          ART\_DIRECTOR);
    \item PATCH /api/v1/events/\{id\}~--- частичное обновление
          полей мероприятия;
    \item DELETE /api/v1/events/\{id\}~--- удалить мероприятие
          (только ART\_DIRECTOR).
\end{itemize}

На стороне фронтенда список мероприятий отображается в виде таблицы
с колонками «Название», «Дата», «Категория», «Ответственный» и
«Статус». Кнопки «Изменить» и «Удалить» отображаются только у тех
строк, к которым у текущего пользователя есть соответствующие права.


Для каждого мероприятия предусмотрена возможность прикрепления
произвольных файлов (документов, изображений и т.д.). Модель
File хранит путь до файла на диске, оригинальное имя,
MIME-тип и временную метку загрузки. Загруженные файлы сохраняются
в директории uploads/\{event\_id\}/ с уникальным именем
на основе UUID, что исключает конфликты имён.

REST API работы с файлами (app/api/files.py) включает:
\begin{itemize}
    \item POST /api/v1/events/\{id\}/files~--- загрузка файла
          (доступна ответственному сотруднику или арт-директору);
    \item GET /api/v1/files/\{id\}/download~--- скачивание
          файла по идентификатору;
    \item DELETE /api/v1/files/\{id\}~--- удаление файла
          (физическое удаление с диска и из базы данных).
\end{itemize}

На странице карточки мероприятия файлы отображаются в виде таблицы,
где каждое имя является ссылкой для скачивания. Загрузка новых файлов
выполняется асинхронно через Fetch API без перезагрузки
страницы (листинг~\ref{lst:upload_js}).

\begin{lstlisting}[caption={JavaScript-функция загрузки файлов}, label={lst:upload_js}]
async function uploadFiles() {
    const input = document.getElementById('file-upload');
    if (!input || !input.files.length) {
        alert('Выберите файлы'); return;
    }
    let hasError = false;
    for (const file of input.files) {
        const fd = new FormData();
        fd.append('file', file);
        const r = await fetch(
            '/api/events/{{ event.id }}/files',
            { method: 'POST', body: fd }
        );
        if (!r.ok) hasError = true;
    }
    if (hasError) alert('Не удалось загрузить файлы');
    location.reload();
}
\end{lstlisting}
\newpage
\subsection{Реализация механизма формирования и выгрузки отчётов в формате Excel}

В соответствии с задачей~3.4 реализован механизм генерации отчётов
в формате Excel (.xlsx). Для работы с табличными
данными используется библиотека pandas, а для записи
файлов~--- openpyxl.

Все конечные точки отчётности расположены в модуле
app/api/reports.py. Ответ каждого эндпоинта является
потоковым (StreamingResponse) с правильно установленными
заголовками Content-Disposition и Content-Type,
что обеспечивает автоматическое скачивание файла браузером
(листинг~\ref{lst:xlsx_response}).

\begin{lstlisting}[caption={Вспомогательная функция формирования Excel-ответа}, label={lst:xlsx_response}]
def _xlsx_response(buf: io.BytesIO, filename: str) -> StreamingResponse:
    buf.seek(0)
    return StreamingResponse(
        buf,
        media_type=(
            "application/vnd.openxmlformats-"
            "officedocument.spreadsheetml.sheet"
        ),
        headers={
            "Content-Disposition":
                f'attachment; filename="{filename}"'
        },
    )
\end{lstlisting}


Эндпоинт GET /api/v1/reports/monthly формирует сводку за
текущий календарный месяц. С помощью функций extract
SQLAlchemy выполняется выборка мероприятий, чья дата проведения
попадает в нужный период. Результирующий файл содержит два листа:
\begin{itemize}
    \item «Итог (месяц год)»~--- сводная таблица: количество
          проведённых и запланированных мероприятий, самая частая
          категория;
    \item «По категориям»~--- детализация проведённых мероприятий
          в разрезе категорий.
\end{itemize}

Эндпоинт GET /api/v1/reports/employees агрегирует данные
по всем сотрудникам (роли, отличные от ADMIN). Для каждого
сотрудника рассчитываются:
\begin{itemize}
    \item количество созданных мероприятий;
    \item количество проведённых мероприятий (в роли ответственного);
    \item доля отклонённых мероприятий от общего числа созданных, \%;
    \item доля незавершённых мероприятий (статус PLANNED) от
          суммы запланированных и проведённых, \%.
\end{itemize}

Пример вычисления процента отклонённых мероприятий приведён в
листинге~\ref{lst:rejection_rate}.

\begin{lstlisting}[caption={Вычисление доли отклонённых мероприятий}, label={lst:rejection_rate}]
total_by_author = (
    df_events.groupby("author_id").size()
    .reset_index(name="total")
    .rename(columns={"author_id": "id"})
)
rejected_by_author = (
    df_events[df_events["status"] == rejected_str]
    .groupby("author_id").size()
    .reset_index(name="rejected")
    .rename(columns={"author_id": "id"})
)
rej = total_by_author.merge(
    rejected_by_author, on="id", how="left"
).fillna(0)
rej["% отклоненных"] = (
    rej["rejected"] / rej["total"] * 100
).round(1)
\end{lstlisting}

Файл также содержит второй лист «По категориям» с разбивкой
проведённых мероприятий каждого сотрудника по категориям в виде
сводной (pivot) таблицы.

Эндпоинт GET /api/v1/reports/structure-units формирует
сводку в разрезе структурных подразделений (СП). Для каждого
подразделения определяются:
\begin{itemize}
    \item количество сотрудников;
    \item суммарное количество проведённых мероприятий, ответственный
          за которые числится в данном СП;
    \item разбивка проведённых мероприятий по категориям.
\end{itemize}

\newpage
\subsection{Разработка системы мониторинга статусов мероприятий}

Согласно задаче~3.5, в системе реализован жизненный цикл мероприятия
на основе конечного автомата (State Machine). Допустимые переходы между
статусами описаны в виде словаря \_TRANSITIONS в модуле
app/core/permissions.py и представлены на
схеме~\ref{lst:transitions}.

\begin{lstlisting}[caption={Таблица допустимых переходов статусов}, label={lst:transitions}]
_TRANSITIONS: dict[EventStatus, set[EventStatus]] = {
    EventStatus.ON_APPROVAL: {
        EventStatus.APPROVED, EventStatus.REJECTED
    },
    EventStatus.APPROVED:  {EventStatus.PLANNED},
    EventStatus.PLANNED:   {EventStatus.COMPLETED},
    EventStatus.COMPLETED: set(),
    EventStatus.REJECTED:  set(),
}
\end{lstlisting}

Перевод мероприятия в статусы APPROVED и REJECTED
доступен исключительно художественному руководителю
(ART\_DIRECTOR). Перевод из PLANNED в
COMPLETED выполняет ответственный сотрудник или художественный руководитель.

При каждой смене статуса в таблицу event\_history автоматически
добавляется запись с указанием пользователя, совершившего изменение,
предыдущего и нового статуса, а также необязательного комментария.

На странице карточки мероприятия отображается:
\begin{itemize}
    \item текущий статус в виде цветного значка-бейджа;
    \item выпадающий список доступных переходов (формируется функцией
          allowed\_statuses на основе роли пользователя и
          текущего статуса);
    \item поле для комментария при смене статуса;
    \item таблица истории изменений со временем, автором и
          комментарием каждого перехода.
\end{itemize}